\subsection{Assignment 5 -- Data Preparation and Dataset Splitting}

L’Assignment 5 richiedeva lo sviluppo di un programma Python in grado di
suddividere i dati di input in file distinti, uno per ciascuna dimensione
e per la fact table del Data Warehouse progettato nell’Assignment~4,
al fine di preparare i dataset per le successive fasi di caricamento.

Per soddisfare tale richiesta, il processo di \textit{data preparation}
è stato articolato in due fasi consecutive e complementari:
\begin{itemize}
    \item pulizia e normalizzazione dei dati;
    \item suddivisione ed esportazione dei dataset, uno per ciascuna tabella
    del Data Warehouse.
\end{itemize}

\subsubsection*{Pulizia e normalizzazione dei dati}

La prima fase è stata implementata tramite lo script
\texttt{prepare\_dw\_datasets.py}, il cui obiettivo è rendere i dati coerenti
con lo schema multidimensionale progettato.

Lo script carica i dataset in formato JSON dalla directory
\texttt{dataset/correct\_ids} ed esegue una serie di operazioni di pulizia
e normalizzazione prima dello \textit{split} vero e proprio.
In particolare, vengono effettuate le seguenti operazioni:
\begin{itemize}
    \item uniformazione dei tipi di dato (interi, float e stringhe);
    \item gestione consistente dei valori mancanti e non validi;
    \item normalizzazione degli attributi testuali;
    \item correzione di incongruenze sintattiche ereditate dai dataset originali;
    \item generazione e validazione delle chiavi surrogate, in linea con lo schema
    del Data Warehouse;
    \item normalizzazione della dimensione temporale e derivazione
    dell’attributo \textit{Season}.
\end{itemize}

Il risultato di questa fase è un insieme di dataset puliti e coerenti,
salvati nella directory \texttt{dataset/cleaned\_json}, che rappresentano
una versione intermedia pronta per essere suddivisa secondo le tabelle
del Data Warehouse.

\subsubsection*{Esportazione}

La seconda fase è stata implementata tramite lo script
\texttt{export\_dw\_csv.py}, che realizza esplicitamente quanto richiesto
dalla consegna, ovvero la suddivisione dei dati in file separati,
uno per ciascuna tabella del Data Warehouse.

A partire dai dataset normalizzati, lo script genera un file CSV per ogni
dimensione e per la fact table. Ogni file CSV contiene esclusivamente
gli attributi previsti dallo schema del Data Warehouse ed è costruito
in modo da rispettarne la granularità.

In particolare, la fact table \textit{FactParticipation} mantiene la
granularità definita nell’Assignment~4, ovvero:
\begin{quote}
\emph{una riga per partecipazione artista--canzone}.
\end{quote}
Essa include la misura \textit{Streams1Month} e l’attributo \textit{IsPrimary},
modellato come \textit{degenerate dimension}.

I file CSV prodotti costituiscono quindi l’output finale dell’Assignment~5
e rappresentano l’input diretto per le successive fasi di caricamento
del Data Warehouse.

L’Assignment 5 ha permesso di trasformare i dati di partenza in un insieme
di file strutturati, uno per ciascuna tabella del Data Warehouse,
come richiesto dalla consegna. La separazione tra fase di normalizzazione
e fase di esportazione consente di ottenere dati puliti, coerenti e
direttamente utilizzabili, facilitando le successive attività di
popolamento e analisi.
